% !TEX root = ../notes.tex

\documentclass[../notes.tex]{subfiles}

\begin{document}

\section{March 30}
Here we go.

\subsection{Forms of Semisimple Lie Algebras}
We will work over a field $K$ of characteristic zero. Recall that we already have a construction of many semisimple Lie algebras.
\begin{definition}[split form]
	Fix a field $K$ and a Dynkin diagram $D$. Then the Serre presentation defines a finite-dimensional Lie algebra $\mf g_D$ with Dynkin diagram $D$ over $K$. We call this the \textit{split form}.
\end{definition}
\begin{example}
	Over algebraically closed fields, we showed that all Lie algebras with Dynkin diagram $D$ are split forms.
\end{example}
\begin{example}
	There are also non-split forms: for example, $\mf{su}_2$ is a real Lie algebra with the same Dynkin diagram as $\mf{sl}_2$, but these are not isomorphic real Lie algebras.
\end{example}
Anyway, the point is that we may classify Lie algebras $\mf a$ over a field $K$ of characteristic zero by classifying the Lie algebras $\mf a\otimes_K\overline K$ over $\overline K$ (which is split!) and then classifying the Lie algebras that base-change to the split form. It will be useful to have some language for this situation.
\begin{definition}[form]
	Fix a Galois extension $L/K$ of fields, and choose a Dynkin diagram $D$. Then a Lie algebra $\mf g$ over $K$ \textit{splits} over $L$ if and only if $\mf g_L$ is isomorphic to the split form. We say that $\mf g$ is a \textit{form} of $\mf g_D$.
\end{definition}
\begin{remark}
	Because all forms split over $\overline K$, one can show that all forms split over a finite Galois extension.
\end{remark}
It will turn out that the Lie algebras $\mf a$ over $K$ with $\mf a\otimes_KL=\mf g_D$ are given by $\mathrm H^1(\op{Gal}(L/K);\op{Aut}\mf g_D)$. Let's recall what $\mathrm H^1$ means for our purposes.
\begin{definition}
	Fix a group $\Gamma$ acting by automorphisms on another group $G$. Then we define $\mathrm H^1(\Gamma;G)$ to classify homomorphisms $\Gamma\to\Gamma\rtimes G$ for which the composite back down to $\Gamma$ is the identity, up to suitable conjugacy.
\end{definition}
\begin{remark}
	Let's explain how $1$-cocycles enter the picture. Such a homomorphism $\varphi\colon\Gamma\to\Gamma\rtimes G$ necessarily has $\varphi(\gamma)=(\eta(\gamma),\gamma)$, where $\eta\colon\Gamma\to G$ is a map which now must satisfy
	\[(\eta(\gamma_1\gamma_2),\gamma_1\gamma_2)=(\eta(\gamma_1),\gamma_1)(\eta(\gamma_2),\gamma_2),\]
	which is equivalent to $\eta(\gamma_1\gamma_2)=\eta(\gamma_1)\cdot\gamma_1(\eta(\gamma_2))$.
\end{remark}
\begin{remark}
	Let's explain how $1$-coboundaries enter the picture. Well, we are conjugating a map $\gamma\mapsto(\eta(\gamma),\gamma)$ by an element $g\in G$, which corresponds to the map $\gamma\mapsto\left(g\eta(\gamma)\gamma(g)^{-1},\gamma\right)$. Thus, conjugacy of homomorphisms corresponds to two $1$-cocycles being the same up to $1$-coboundary.
\end{remark}
Let's explain how to relate forms to cohomology. The idea is to note that the Galois group acts on the base-change.
\begin{definition}[twisted linear]
	Fix a Galois extension $L/K$ of fields, and let $\Gamma$ be the Galois group. Then a \textit{twisted linear action} of $\Gamma$ on a vector space $V$ over $L$ is one which is additive, and
	\[\gamma(\ell x)=\gamma(\ell)\gamma(x)\]
	for any $\gamma\in\Gamma$ and $\ell\in L$ and $x\in V$.
\end{definition}
\begin{example}
	For any vector space $V$ over $K$, there is a trivial twisted linear action on $V\otimes_KL$ given by having $\Gamma$ act directly on the $L$-component. By computing a basis, one finds that $V=(V\otimes_KL)^\Gamma$.
\end{example}
\begin{proposition} \label{prop:classify-ss-forms}
	Fix a Galois extension $L/K$ of fields, and choose a Dynkin diagram $D$. Then the pointed set $\mathrm H^1(\op{Gal}(L/K);\op{Aut}\mf g_D)$ classifies Lie algebras $\mf g$ over $K$ with Dynkin diagram $D$ which split over $L$.
\end{proposition}
\begin{proof}
	Set $\Gamma\coloneqq\op{Gal}(L/K)$. For any Lie algebra $\mf g$ over $K$, we see that $\Gamma$ acts on $\mf g_L$ via a twisted linear action, and then one can recover $\mf g$ from $\mf g_L$ by taking invariants. Conversely, any twisted linear action of $\Gamma$ on a Lie algebra $\mf g_L$ over $L$ can take $\Gamma$-invariants to recover a Lie algebra over $K$.

	Now, fix a Lie algebra $\mf g\coloneqq\mf g_D$ over $K$, and let $\rho$ be a twisted linear action of $\Gamma$ on $\mf g_L$. Further, let $\rho_0$ be the trivial twisted linear action. Then there is $\eta\colon\Gamma\to\op{Aut}(\mf g)$ for which $\rho(g)=\eta(g)\rho_0(g)$ for all $g\in\Gamma$. Note that $\rho$ is a homomorphism if and only if $\rho(gh)=\rho(g)\rho(h)$ for all $g$ and $h$, which we can see to be equivalent to having
	\[\eta(gh)=\eta(g)\circ g\eta(h),\]
	where $g\eta(h)\coloneqq\rho_0(g)\eta(h)\rho_0(g)^{-1}$. Thus, $\eta$ is a $1$-cocycle in $\mathrm Z^1(\Gamma;\op{Aut}(\mf g))$.
	
	Conversely, $\eta_1$ and $\eta_2$ define isomorphic $\rho_1$ and $\rho_2$ if and only if there is $a\in\op{Aut}(\mf g)$ for which $\rho_1(g)\circ a=a\circ\rho_2(g)$ for all $g$. One can expand that this is equivalent to having $\eta_1(g)=a\eta_2(g)ga^{-1}$ for all $g$, which is equivalent to $\eta_1$ and $\eta_2$ being cohomologous. In total, we have obtained a bijection from isomorphism classes of Lie algebras $\mf g'$ over $D$ with $\mf g'_L=\mf g_L$ and $\mathrm H^1(\Gamma;\op{Aut}\mf g)$, as required.
\end{proof}
\begin{remark}
	The same argument works for any linear algebraic structure. Namely, one can replace Lie algebras with vector spaces, quadratic spaces, and so on.
\end{remark}

\subsection{Real Forms}
In the sequel, we will focus on the Galois extension $\CC/\RR$. Let's work out \Cref{prop:classify-ss-forms} in this case.
\begin{example}
	Consider the Galois extension $\CC/\RR$. In this case, \Cref{prop:classify-ss-forms} (combined with \Cref{prop:aut-lie-alg}) tells us that real forms of a complex semisimple Lie algebra $\mf g_D$ are given by
	\[\mathrm H^1(\{\pm1\};\op{Aut}(D)\ltimes G_{\mathrm{ad}}),\]
	where $-1\in\{\pm1\}$ acts by complex conjugation. Notably, a $1$-cocycle $\eta\colon\{\pm1\}\to\op{Aut}(\mf g_D)$ must have $\eta(1)=\id$, so we only have to worry about the data of $s\coloneqq\eta(-1)$. The cocycle condition amounts to requiring $s\overline s=1$. Further, note that we only care about $\eta$ up to conjugation, which means that we are only interested in $s$ up to the conjugation $s\mapsto as\overline a^{-1}$, where $a\in\op{Aut}(\mf g_D)$. Thus, real forms of $G$ are in bijection with $s\in\op{Aut}(D)\ltimes G_{\mathrm{ad}}$ for which $s\overline s=1$, considered up to twisted conjugation.
\end{example}
\begin{notation}
	We work over $\RR$. Fix a Dynkin diagram $D$. For $s\in\op{Aut}(\mf g_{D,\CC})$ for which $s\overline s=1$, we let $\mf g_s$ denote the corresponding real form.
\end{notation}
This $s$ is a $\CC$-linear transformation, but it will turn out to be convenient to instead consider ``antilinear'' transformations.
\begin{definition}[antilinear]
	Fix a complex vector space $V$. Then an $\mathbb R$-linear map $\sigma\colon V\to V$ is \textit{antilinear} if and only if $\sigma(zv)=\overline z\sigma(v)$ for $z\in\CC$.
\end{definition}
\begin{example}
	Suppose that $V$ is a real vector space. Then a linear map $s\colon V_\RR\to V_\RR$ induces an antilinear map $\sigma_s\colon V_\CC\to V_\CC$ given by composing with complex conjugation.
\end{example}
\begin{remark}
	By definition, we have that $\mf g_s$ consists of those $X\in\mf g_{D,\CC}$ for which $\overline X=s(X)$, which is equivalent to $X\in\mf g^{\sigma_s}$.
\end{remark}
Let's now compute some real forms. Note that complex conjugation acts trivially on the Dynkin diagram itself, so our $s\in\op{Aut}(\mf g_{D,\CC})$ induces an automorphism $s_0\in\op{Aut}(D)$ with $s_0^2=1$. It is not too hard to classify such involutions of the Dynkin diagram. Thus, a semisimple Lie algebra $\mf g$ over $\RR$ will have $\mf g_\CC$ be given by some Dynkin diagram $D$ with an automorphism $s_0$; by organizing the Dynkin diagram into $s_0$-invariant pieces, we see that $\mf g$ is a sum of simple Lie algebras.

Further, if $\mf g$ is actually simple, then there are two cases.
\begin{example}
	It is possible for $\mf g_\CC$ to consists of two copies of a connected Dynkin diagram $D_0$, meaning that $D=D_0\sqcup D_0$, and $s_0$ swaps the two pieces. Thus, we see that $G_{\mathrm{ad}}=K_{\mathrm{ad}}\times K_{\mathrm{ad}}$, where $s_0$ swaps the two factors. Let's now write out the condition $s\overline s=1$. We know that $s=s_0(g,h)$ where $g,h\in K_{\mathrm{ad}}$. Now, $s\overline s=1$, we see $s_0(g,h)=s_0(\overline g\overline h)=1$, which means that $h=\overline g^{-1}$. But then $s=s_0\left(g,\overline g^{-1}\right)=(1,g)s_0\overline{(1,g)}^{-1}$, so $s$ is cohomologous with $s_0$!

	Thus, we see that we only have a single real form here. Taking $s_0$-invariants, we see that our Lie algebra $\mf g_{s_0}$ is simply given by $\mf k=\op{Lie}K$. Explicitly, we have taken some simple complex Lie algebra $\mf k$ and forgotten the complex structure to produce a simple real Lie algebra.
\end{example}
Alternatively, we could have $\mf g_\CC$ be a simple complex Lie algebra, and $s_0$ acts nontrivially on the Dynkin diagram. This turns out to require some interesting combinatorics.

\subsection{Inner Forms}
Let's start by organizing our forms by their image in $\op{Aut}D$.
\begin{definition}[inner form]
	Fix a Dynkin diagram $D$ and a field $K$. A form $\mf g_s$ with Dynkin diagram $D$ is \textit{inner} if and only if the image of $s$ in $\op{Aut}(D)$ is trivial.
\end{definition}
\begin{remark}
	It is equivalent to ask for $s$ to be in $G_{\mathrm{ad}}$.
\end{remark}
\begin{definition}[inner class]
	Fix a Dynkin diagram $D$ and a field $K$. Two forms $\mf g_s$ and $\mf g_{s'}$ with Dynkin diagram $D$ are in the same \textit{inner class} if and only if $s$ and $s'$ have the same image in $\op{Aut}(D)$.
\end{definition}
\begin{remark}
	Thus, the inner forms make up the collection of those forms in the same inner class as the split form.
\end{remark}
Each inner class admits a distinguished element.
\begin{definition}[quasisplit]
	Fix a Dynkin diagram $D$ and a field $K$. A form $\mf g_s$ is \textit{quasisplit} if and only if the image of $s$ in $G_{\mathrm{ad}}$ is trivial (up to conjuacy).
\end{definition}
\begin{remark}
	Each inner class consists of elements of the form $(s_0,g)\in \op{Aut}(\mf g)\rtimes G_{\mathrm{ad}}$, where $s_0$ is fixed. Thus, each inner class admits a unique quasisplit form.
\end{remark}
The above discussion works for general fields.

\subsection{The Compact Form}
We now define a special involution of complex Lie algebras.
\begin{definition}[Cartan involution]
	Fix a semisimple Lie algebra $\mf g$ over $\CC$. Then we define the \textit{Cartan involution} $\tau\colon\mf g\to\mf g$ to be given on the Serre presentation
	\[\begin{cases}
		\tau(h_j)\coloneqq-h_j, \\
		\tau(e_j)\coloneqq-f_j, \\
		\tau(f_j)\coloneqq-e_j
	\end{cases}\]
	for each $j$.
\end{definition}
\begin{remark}
	We will not bother to check that $\tau$ preserves the Serre relations, which holds by their symmetry. Note that $\tau\overline\tau=1$.
\end{remark}
With an involution, we get a real form.
\begin{definition}[compact form]
	Fix a semisimple Lie algebra $\mf g$ over $\CC$. Then the \textit{compact real form} $\mf g^c$ is the real form $\mf g_\tau$.
\end{definition}
It will turn out that $\mf g^c$ corresponds to a compact group, which explains the language. We will prove this shortly, but let's start with an example.
\begin{example} \label{ex:compact-sl2}
	Consider $\mf g=\mf{sl}_{2}(\CC)$. Then $\tau(h)=-h$ and $\tau(e)=-f$ and $\tau(f)=-e$, so $\mf g^c=\{X\in\mf g:\overline X=\tau(X)\}$ admits a basis over $\RR$ given by the matrices
	\[\left\{\begin{bmatrix}
		i \\ & -i
	\end{bmatrix},\begin{bmatrix}
		& 1 \\ -1
	\end{bmatrix},\begin{bmatrix}
		& i \\ i
	\end{bmatrix}\right\}.\]
	Thus, we see that $\mf{sl}_2(\CC)^c=\mf{su}(2)$.
\end{example}
\begin{theorem} \label{thm:compact-form-compact}
	Fix a complex semisimple Lie algebra $\mf g$. Then the Killing form of $\mf g^c$ is negative definite.
\end{theorem}
\begin{proof}
	The root decomposition allows us to write
	\[\mf g=\mf h\oplus\bigoplus_{\alpha\in\Phi_+}(\CC e_\alpha\oplus\CC f_\alpha).\]
	Now, $\tau$ acts on each of the above summands, so we find that
	\[\mf g^c=\mf h^c\oplus\bigoplus_{\alpha\in\Phi_+}(\CC e_\alpha\oplus\CC f_\alpha)^c.\]
	Now, this is an orthogonal decomposition for the Killing form (indeed, this is already true for the root decomposition), so it is enough to show the negative-definiteness on each summand.
	\begin{itemize}
		\item We show that the Killing form is negative-definite on $\mf h^c$. Here, $\tau$ acts by $-1$ on $\mf h$, we see that $\mf h^c$ admits a basis given by the $ih_\bullet$s. On this basis, we choose some $i\sum_jt_jh_j\in\mf h^c$, and then we can compute the Killing form by
		\[\tr_{\mf g}\left(i\sum_jt_jh_j\right)^2=-\tr_{\mf g}\left(\sum_jt_jh_j\right)^2.\]
		Now, via the root decomposition, we may compute this trace as $-\sum_\alpha\left(\sum_jt_j\alpha(h_j)\right)^2$, which is negative.

		\item We show that the Killing form is negative-definite on $(\CC e_\alpha\oplus\CC f_\alpha)^c$ when $\alpha$ is a simple root $\alpha_j$. Then by definition of $\tau$, we see that $\mf g^c\cap\mf{sl}_{2j}$ is isomorphic to $\mf{su}(2)$, as we just computed in \Cref{ex:compact-sl2}. Here, the Killing form takes the form
		\[K(X,Y)=\tr_{\mf g}(\rho_{\mf g}(X),\rho_{\mf g}(Y))\]
		for $X,Y\in\mf{su}(2)$. But all representations of $\mf{su}(2)$ lift to the compact group $\op{SU}(2)$ and thus are unitary, so $\rho_{\mf g}$ is skew-Hermitian and hence have imaginary eigenvalues, so $\tr\rho_{\mf g}(X)^2<0$ for nonzero $X$, so the Killing form is negative-definite!
		
		\item It remains to handle the general case for the spaces $(\CC e_\alpha\oplus\CC f_\alpha)^c$. Well, recall the construction of $w_j$ from \Cref{prop:compute-weyl-group}. Note that $\begin{bsmallmatrix}
			0 & 1 \\ -1 & 0
		\end{bsmallmatrix}$ lives in $\op{SU}(2)$. Now, any element $G_{\mathrm{ad}}$ preserves the Killing form on its Lie algebra $\mf g$, so $w_j$ preserves the Killing form of $\mf g^c$. It follows that any $\widetilde w$ preserves the Killing form. The point is that any root $\alpha$ can be permuted with $W$ to be simple, and we now know that the Weyl group action preserves the Killing form, so we reduce to the previous case!
		\qedhere
	\end{itemize}
\end{proof}
We may now explain how $\mf g^c$ corresponds to a compact real group.
\begin{definition}
	Fix a complex semisimple Lie algebra $\mf g$. Then we define $G_{\mathrm{ad}}^c$ to be the real Lie subgroup of $G_{\mathrm{ad}}$ corresponding to the Lie algebra $\mf g^c$.
\end{definition}
\begin{remark}
	By construction, $G_{\mathrm{ad}}^c=\{X\in G_{\mathrm{ad}}:\tau(X)=\overline X\}^\circ$, where the action of $\tau$ is by conjugation. Indeed, one shows this by taking the differential of the relation $\tau(X)=\overline X$. In fact, $\{X\in G_{\mathrm{ad}}:\tau(X)=\overline X\}$ is connected, which we will show later.
\end{remark}
\begin{remark}
	Note that $G_{\mathrm{ad}}^c$ acts by negative-definite automorphisms on $\mf g$, so $G_{\mathrm{ad}}^c$ embeds into an orthogonal group, so it is compact.
\end{remark}
\begin{remark}
	Later, we will show that the universal cover of $G_{\mathrm{ad}}^c$ continues to be compact. Thus, its representation theory is semisimple by \Cref{cor:reps-semisimple}, so the representation theory of $\mf g^c$ is semisimple, so the representation theory of $\mf g$ is semisimple!
\end{remark}
One can compute this for the classical types without too much effort. We will just give the answers.
\begin{example}
	If $\mf g=\mf{sl}_n(\CC)$, then $G_{\mathrm{ad}}^c$ turns out to be $\op{PSU}(n)$.
\end{example}
\begin{example}
	If $\mf g=\mf{so}_n(\CC)$, then $G_{\mathrm{ad}}^c$ is $\op{SO}_n(\CC)$ if $n$ is odd, and it is $\op{SO}_n(\CC)/\{\pm1\}$ if $n$ is even.
\end{example}
\begin{example}
	Lastly, if $\mf g=\mf{sp}_{2n}(\CC)$, then $G_{\mathrm{ad}}^c$ is $(\op{Sp}_{2n}(\CC)\cap\op U(2n))/\{\pm1\}$. Alternatively, one can describe these as the unitary $n\times n$ matrices with coefficients in the quaternions.
\end{example}
\begin{remark}
	The compact form is not quasisplit. Indeed, the quasisplit real form has nilpotent elements (it comes from an automorphism of the Dynkin diagram, so it has the element $e$ from the principal $\mf{sl}_2$-triple), so the corresponding real Lie group is not compact.
\end{remark}

\end{document}